%% kapa: exam=1 round=23 year=2012 period=2 form=A subject=accounting subject_ko=회계학 seq=28 number=28
%% source: Q-Net 감정평가사 기출문제 (cst00309) artlSeq=5052580 file=raw/1008/5052580/2012년 제23회 감정평가사 제1차시험 기출문제(2교시 A형).pdf
%% conversion: kordoc 4.24.51 -> tools/gen1.py (auto); review=proofread-sonnet; flags=none
\begin{question}{28}
(주)강북은 20×1년 1월 1일 건물(매입가격 ￦10,000,000, 매입부대비용 ￦150,000)을 취득하는 과정에서 공채(액면금액 ￦100,000, 이자율 7\%, 만기 3년, 매년말 이자지급 조건)를 불가피하게 액면금액에 인수하였다. 공채 취득시점의 (주)강북에 적용될 시장이자율은 10\%였고, 건물대금과 공채대금은 모두 현금으로 지급하였다. 취득한 공채는 만기보유금융자산으로 분류되었고, 건물은 정액법으로 감가상각(내용연수 5년, 잔존가치 없음)한다. 이의 회계처리에 관한 설명으로 옳은 것은? (단, 계산금액은 소수점 첫째자리에서 반올림한다. 기간말 단일금액의 현가계수(3기간, 10\%)는 0.7513, 정상연금의 현가계수(3기간, 10\%)는 2.4868이다.)
\begin{choices}
  \item 20×1년 12월 31일 인식할 건물의 감가상각비는 ￦2,030,000이다.
  \item 20×1년 1월 1일 만기보유금융자산의 장부금액은 ￦93,500이다.
  \item 20×1년 12월 31일 인식될 이자수익은 ￦7,000이다.
  \item 20×1년 1월 1일 건물의 취득원가는 ￦10,157,462이다.
  \item 20×1년 12월 31일 만기보유금융자산의 장부금액은 ￦95,000이다.
\end{choices}
\end{question}
\answer{4}
